% ══════════════════════════════════════════════════════════════════════════════
% CHAPTER 8 — CONCLUSION  [v31 — revised for defence]
% ══════════════════════════════════════════════════════════════════════════════

\section{Summary of findings}
\label{sec:conc_summary}

This thesis studied whether the hawkish--dovish tone of central bank
communication contains information for financial markets and for future policy
outcomes. The direct answer to RQ1 is yes, but only conditionally: text explains
some announcement-window market reactions after controlling for rate surprises,
with the clearest evidence for FOMC two-year yields and ECB equities rather than
for all assets. The direct answer to RQ2 is positive for policy decisions and
negative for yields: language helps forecast next-meeting and two-meeting-ahead
rate decisions in an expanding-window setting, but all sixteen yield forecasts
(two methods, four corpora, two horizons) deliver negative direct out-of-sample
$R^2$.

The empirical analysis used a corpus of 884 FOMC and ECB statements and press
conference transcripts covering 1994--2026. Three text-scoring methods were
applied to the same documents: a transparent dictionary score, a class-balanced
Random Forest score based on Lasso-selected bigrams, and a BERT
sentence-level score aggregated to the document level. The scores were then
compared through agreement tests, event-study regressions, method horse races,
decomposition exercises, and expanding-window forecasts. Because the analysis
spans four corpora, multiple outcomes, three methods, and several regression
variants---generating in excess of 200 individual specifications---individual
significant coefficients should be interpreted cautiously, and the emphasis
throughout is on patterns that are consistent across methods and outcomes.

The first research question asked whether central bank tone explains
announcement-window financial market reactions after controlling for the
contemporaneous monetary policy surprise. The answer is concentrated rather
than universal. For FOMC statements, the strongest evidence is at the short end
of the Treasury curve. The two-year yield regressions show that Random Forest
and BERT tone are positively associated with announcement-window yield
changes (adjusted $R^2 = 0.249$ in the full horse race), and this result is
robust to Newey-West standard errors and fixed-effect controls. The ten-year
yield evidence is weaker. This maturity pattern is economically consistent with
the view that policy communication matters most for assets tied closely to the
expected near-term policy path.

For ECB statements, the strongest market-reaction result is in equities rather
than yields. The Random Forest score is positively associated with Euro Stoxx 50
responses ($p=0.028$), and this result survives the three-way horse race
($p=0.045$). The positive sign is consistent with an information-effect
interpretation, though the result becomes marginal ($p=0.102$) when chair and
era fixed effects are added, and the two channels cannot be distinguished
structurally in the present framework. Two other ECB outcomes---OIS 10Y and
EUR/USD---have negative adjusted $R^2$ in all specifications, confirming that
text adds no explanatory power for these assets after adjusting for degrees of
freedom. The forward/current split provides partial evidence for the
dual-channel view: forward-looking hawkish language is negative for equities,
while hawkish current-condition language is positive, and this pattern survives
fixed-effect controls.

The press conference evidence is richer but less stable. Press conferences
contain longer, less formulaic, and more explanatory text than statements, which
makes them valuable for semantic methods such as BERT. However, they
also combine policy explanation, risk assessment, macroeconomic interpretation,
and answers to journalist questions. In the FOMC press-conference sample
($N=91$), the Random Forest score is significantly associated with Treasury
yields, but the signs are negative---the opposite of the standard prediction.
This result is robust to outlier removal and Newey-West standard errors, and
it illustrates why the mapping from tone to market reaction is more complex
than the methodology's theoretical framework suggests.

The second research question asked whether central bank language predicts
near-term policy outcomes and two-year yield changes in a strictly
expanding-window setting. The answer is unambiguous for rates and clearly
negative for yields. Both Random Forest and BERT forecasts deliver
positive direct out-of-sample $R^2$ values for next-meeting and two-meeting-ahead
rate decisions in all eight corpus--method pairs. BERT generally performs
best for these rate-decision targets, especially in press conferences. However,
all three text scores are highly autocorrelated ($\rho_1 \approx 0.55$--$0.87$),
so part of the forecast power may reflect the text score acting as a slow-moving
proxy for the policy stance rather than capturing meeting-specific forward
guidance. The incremental $R^2$ analysis, which controls for the current rate
decision, partially addresses this concern and confirms that the text forecast
adds significant information beyond inertia alone. By contrast, direct
out-of-sample $R^2$ values for two-year yield changes are negative in all
sixteen cases. Communication helps forecast the central bank's own subsequent
policy choices, but it does not forecast the market's future repricing.

The method-comparison results are also central to the thesis. The dictionary,
Random Forest, and BERT scores are related, but they do not measure the
same object. The dictionary captures explicit stance-bearing vocabulary and is
the most transparent method. The Random Forest captures language that is
predictive of the contemporaneous decision sign, which explains its high
agreement with rate decisions but also limits its interpretation as an
independent validation device; moreover, approximately 18--20\% of its feature
importance in FOMC and ECB statements is attributable to leakage-adjacent terms
that correlate with speakers or eras rather than economic stance. BERT
captures sentence-level semantic information and is particularly useful in the
expanding-window rate-decision forecasting exercise, though its systematic
dovish shift (81 of 91 FOMC press conferences labelled dovish) reflects domain
transfer from the FOMC-trained classifier. No method dominates unconditionally.
The contribution of the three-way design is precisely that it shows when the
measures agree, when they diverge, and how their economic interpretation differs.

\section{Contributions}
\label{sec:conc_contributions}

The thesis makes four main contributions. First, it provides a long-sample,
cross-bank comparison of FOMC and ECB communication using a consistent empirical
framework. The sample covers conventional policy cycles, the global financial
crisis, zero and negative-rate periods, unconventional policy, the post-COVID
tightening cycle, and the recent easing phase. This makes the analysis broader
than a single-institution or single-regime study.

Second, the thesis combines three text-measurement approaches on the same corpus.
Rather than comparing results across studies that differ in sample, data source,
and model design, the dictionary, Random Forest, and BERT measures are
constructed for the same documents and tested in the same regressions. This
allows a direct assessment of the trade-off between transparency, flexibility,
and context sensitivity, while also exposing specific method limitations such as
RF feature leakage and BERT domain transfer.

Third, the thesis decomposes central bank tone into economically interpretable
components. The dictionary separates policy-action, inflation, employment and
real-activity, and residual-risk language. It also separates forward-looking from
current-condition text. These decompositions show that the aggregate
hawkish--dovish score can mask economically important heterogeneity. The ECB
equity result is the clearest example: forward-looking and current-condition
hawkishness carry opposite market signals. For ECB OIS 10Y, the decomposition
turns a model with negative adjusted $R^2$ into one with modest positive
explanatory power (partial $F$-test $p=0.027$).

Fourth, the thesis evaluates predictive content using a strict expanding-window
design. At each forecast date, the model uses only information available before
the forecasted observation. This design is more demanding than in-sample or
out-of-bag evaluation and gives a cleaner answer to whether central bank
language is useful for real-time forecasting. One caveat is that the
BERT sentence classifier is applied once to the full corpus; only the
regression layer that maps the score to forecast targets is fit in an expanding
window. The result is that language predicts future rate decisions but does not
predict subsequent two-year yield changes.

\section{Implications}
\label{sec:conc_implications}

The main methodological implication is that central bank tone should not be
treated as one stable, universal variable. Communication contains policy-action
language, macroeconomic assessment, risk framing, and institutional wording.
These components can move together, but they can also have different market
interpretations. A single net hawkishness score is useful as a summary measure,
but it should be supplemented by decompositions and alternative scoring methods
when the empirical question concerns market reactions.

A related measurement implication is that supervised text classifiers trained on
policy outcomes require careful feature-importance diagnostics. The Random
Forest achieves high in-sample agreement with the rate-decision sign, but
approximately one-fifth of its feature importance is attributable to
leakage-adjacent terms rather than economic stance. Blocklists and
feature-inspection routines are therefore necessary complements to accuracy
metrics, not optional appendix material.

The main economic implication is that markets react not only to whether language
is hawkish or dovish, but also to why it is hawkish or dovish. Hawkish language
based on future tightening can be contractionary for equities. Hawkish language
based on stronger current conditions can be expansionary for equities if markets
interpret it as positive information about the economy. This distinction is
consistent with an information-effect interpretation, though the present
framework cannot structurally separate the policy-news and information-news
channels. It also helps explain why some coefficients---particularly the
negative yield signs in FOMC press conferences---have signs that differ from the
simple policy-tightening prediction.

For practitioners, the results suggest that text-based central bank scores are
most useful when applied narrowly. They should not be read as stand-alone trading
signals across all asset classes. The evidence is strongest for outcomes close
to the communication object: near-term policy decisions and short-maturity yield
responses. It is weaker for longer-maturity yields, foreign exchange, and future
yield changes, where many other macro-financial forces intervene. The practical
value of text models is therefore in complementing, not replacing, high-frequency
surprise measures and institutional judgment.

For central banks, the findings underline that communication is interpreted
through multiple channels. A statement intended to signal vigilance against
inflation may also signal confidence in economic conditions. Similarly, a
press-conference answer intended to preserve flexibility may be interpreted as a
shift in the expected policy path. The market response depends on how investors
separate policy news from information news.

\section{Limitations}
\label{sec:conc_limitations}

The conclusions should be interpreted with several limitations in mind.

All tone measures are constructed variables. The dictionary depends on term
selection and smoothing choices. The Random Forest depends on the decision-sign
training target, feature selection, class balancing, and blocklist. BERT
depends on sentence-level model predictions and the domain on which the model
was trained. Using three methods reduces dependence on any single measurement
choice, but it does not remove measurement error.

The Random Forest classification exercise is descriptive in RQ1. The model is
trained to classify the contemporaneous rate-decision sign, so its high
agreement with rate decisions is a convergence diagnostic rather than independent
evidence that it discovers the true economic stance of communication. Moreover,
approximately 18\% of feature importance in FOMC statements and 20\% in ECB
statements is attributable to leakage-adjacent terms. The \rfscore{} used in
market-reaction regressions is therefore partly influenced by era-specific
vocabulary. The dictionary and BERT scores, which are not trained on the
rate decision, provide a cleaner test of the communication channel.

The analysis reports in excess of 200 individual specifications without a formal
multiple-testing correction. The patterns that are consistent across methods and
corpora---such as the positive rate-decision forecasting result in all eight
corpus--method pairs---are less likely to reflect chance than isolated significant
coefficients, but individual $p$-values should be interpreted conservatively.

The text scores are highly autocorrelated ($\rho_1 \approx 0.55$--$0.87$). Part
of the RQ2 forecasting power may therefore reflect the text score acting as a
slow-moving proxy for the policy stance rather than capturing meeting-specific
forward guidance. The incremental $R^2$ analysis partially addresses this
concern, but a single-lag inertia control may not fully absorb the slow-moving
stance information.

The BERT sentence classifier is applied once to the full corpus and is
not retrained in the expanding-window design. Only the regression layer is fit
on past data. The classifier also exhibits a systematic dovish shift,
particularly in FOMC press conferences, reflecting domain transfer from the
training data.

The ECB policy-rate benchmark is necessarily simplified. A single rate-decision
sign cannot fully capture the ECB's multi-instrument policy environment,
especially after 2014.

Sample sizes are uneven. FOMC press conferences begin only in 2011 ($N=91$),
and the forecast evaluation window is as short as 53 observations for some
targets. ECB long-maturity OIS variables are available for only $N=113$
observations, and the adjusted $R^2$ is negative for these specifications.

The event-study design identifies announcement-window associations rather than
fully structural causal effects. The information-effect interpretation invoked
for the positive ECB equity coefficient is consistent with the data but is not
formally identified; a structural decomposition would be needed to separate the
policy-news and information-news channels.

The direct out-of-sample $R^2$ values are reported as point estimates without
confidence intervals. A Clark-West or bootstrap test would provide a formal
assessment of whether these values are significantly different from zero.

\section{Future research}
\label{sec:conc_future}

Future research could extend the analysis in several directions. The first
priority is to improve the ECB policy benchmark. A multi-instrument ECB measure
that combines the main refinancing rate, deposit facility rate, asset-purchase
announcements, and forward-guidance changes would better match the ECB's
post-2014 policy framework and could improve both agreement tests and
forecasting exercises.

A second extension would be to build institution-specific transformer models.
An ECB-specific sentence classifier would provide a cleaner test of whether
contextual models generalise across central banks or require institution-level
fine-tuning. Retraining the classifier within the expanding-window design would
also address the look-ahead concern.

A third extension would be to rerun the Random Forest with the extended
blocklist identified in the leakage diagnostic. A before-and-after comparison
would quantify how much the 18--20\% leakage-adjacent feature importance
affects the RQ1 regression coefficients and the RQ2 expanding-window forecasts.

Fourth, the out-of-sample $R^2$ values could be tested formally using
Clark-West or Diebold-Mariano tests. Bootstrapping the expanding-window
procedure and reporting confidence bands would strengthen the RQ2 conclusions.

Fifth, the policy-news and information-news channels could be separated
structurally. Combining text scores with sign-restriction or
heteroskedasticity-based identification strategies from the structural
monetary-policy literature would provide a more rigorous test of whether
hawkish language operates through the discount-rate channel, the information
channel, or both.

Sixth, the dictionary could be made dynamic. Era-specific or rolling
dictionaries could test whether the meaning of stance-bearing phrases changes
across conventional policy, the zero lower bound, negative rates,
unconventional policy, and post-COVID tightening.

Seventh, future work could move from document-level to sentence-level market
reaction analysis. Combining sentence-level tone with intraday price movements
would allow a more precise test of which parts of communication move markets.

The results should not be interpreted as showing that text alone causally moves
markets. They show that measured communication tone is associated with selected
announcement-window reactions conditional on rate surprises, and that text-based
forecasts contain information about future policy decisions. The design is
informative about the market relevance and predictive content of communication,
but it does not fully isolate a structural communication shock.

Overall, the thesis shows that central bank language is informative, but not in
a simple or uniform way. Textual tone matters most when the outcome is close to
the communication object itself: future policy decisions and short-maturity
market reactions. Its effects are weaker and more ambiguous when the outcome is
further removed from the announcement, such as long yields, exchange rates, or
future market repricing. The central empirical lesson is therefore not that one
text score can summarise central bank communication, but that different textual
methods reveal different layers of the policy signal---and that each method
carries specific limitations that should be acknowledged alongside its
contributions.